\chapter{能量效率最优化定理证明}\label{app03}
首先考虑$P_T=P_W$且$T_0 \leq \frac{E/N}{P_W}$的情况。根据\eqref{Smax}式，给定数据包块长度$N$的情况下节点生命周期内成功传输数据包数量的期望由下式给出
\begin{equation}
S_{\max}^{N}=\frac{{E}/{N}}{\frac{P_T-P_W}{p_{tw}^N}-\frac{nP_W}{p_{tw}^N\ln{\frac{p_{tw}^N}{1-\epsilon}}}}.
\end{equation}
因为$S_{\max}^{N}(N)$在$N\leq N_{tw}$时随$N$的增加单调递增，且在$N\ge N_{tw}$时随$N$的增加单调递减，那么此时$S$的最大期望在$N=N_{tw}$时取到。令$S_{\max}^{N'}(N)=0$，可以得到
\begin{equation}
nP_W\left[-\frac{\epsilon^{\prime}(N)}{e}\cdot N-\frac{1-\epsilon(N)}{e}\right]=0,
\end{equation}
进而可以得到
\begin{equation}
N\epsilon^{\prime}(N)-\epsilon(N)+1=0.
\end{equation}
解得$N_{tw}$为方程\eqref{N_e}的单实根。

然后考虑$P_T \neq P_W$且$T_0\le\frac{E/N}{P_W-\frac{P_T-P_W}{n}\ln{\frac{p_m^N}{1-\epsilon}}}$的情况。根据\eqref{Smax}，此时的给定数据包块长度$N$的情况下节点生命周期内成功传输数据包数量的期望由下式给出
\begin{equation}\label{eq3}
S_{\max}^{N}=\frac{{E}/{N}}{\frac{P_T-P_W}{p_m^N}-\frac{nP_W}{p_m^N\ln{\frac{p_m^N}{1-\epsilon}}}}.
\end{equation}
因为$S_{\max}^{N}(N)$在$N\leq N_m$时随$N$的增加单调递增，且在$N\ge N_m$时随$N$的增加单调递减，那么此时$S$的最大期望在$N=N_m$时取到。令$S_{\max}^{N'}(N)=0$，可以得到
\begin{equation}\label{eq4}
\begin{aligned}
&(P_T-P_W)\cdot\left(\frac{n-\sqrt{n^2+4n\left(\frac{P_T}{P_W}-1\right)}}{2\left(\frac{P_T}{P_W}-1\right)}\right)^2\cdot[-N\epsilon^{\prime}(N)-(1-\epsilon(N))]+
nP_W\\
&\cdot\frac{n-\sqrt{n^2+4n\left(\frac{P_T}{P_W}-1\right)}}{2\left(\frac{P_T}{P_W}-1\right)}[N\epsilon^{\prime}(N)+(1-\epsilon(N))]=0
\end{aligned}
\end{equation}
进而可以得到
\begin{equation}\label{eq5}
N\epsilon^{\prime}(N)-\epsilon(N)+1=0.
\end{equation}
解得$N_m$为方程\eqref{root1}的单实根。

最后，考虑$P_T \neq P_W$且$\frac{E/N}{P_W-\frac{P_T-P_W}{n}\ln{\frac{p_m^N}{1-\epsilon}}}< T_0\le\frac{{E}/{N}}{P_W}$的情况。此时$S_{\max}^{N}(N)$在$N\leq N_c$时随$N$的增加单调递增，且在$N\ge N_c$时随$N$的增加单调递减，那么此时$S$的最大期望在$N=N_c$时取到。令$S_{\max}^{N'}(N)=0$，可以得到
\begin{equation}\label{eq6}
\left[\frac{nE}{(P_T-P_W)T_0}\cdot\frac{1}{N}-\frac{N\epsilon^{\prime}(N)}{1-\epsilon(N)}-1\right]\left(\frac{E}{NT_0}-P_W\right)=P_W
\end{equation}
解得$N_c$为方程\eqref{root2}的单实根。


综上，定理\ref{Th:Sconstrainedopt}得证。